% =====================================================================================
%  safe_prompt_fields_audit.tex -- audit of safe_prompt_fields.tex (version of 14 October 2026)
%  Build (from this folder):  pdflatex safe_prompt_fields_audit && pdflatex safe_prompt_fields_audit
%  Quotes of the audited note use its own number macros, archived in archive/2026-10-14_pre_audit/tables/.
%  Every recomputed number is a macro of tables/audit_numbers.tex (code/audit_checks.py), and the two
%  evidence tables are tables/audit_probe.tex and tables/audit_attack.tex (same script).
% =====================================================================================
\documentclass[11pt]{article}
\usepackage[margin=2.4cm]{geometry}
\usepackage{amsmath,amssymb,amsthm,mathtools}
\usepackage{bbm}
\usepackage{booktabs}
\usepackage{longtable}
\usepackage{array}
\usepackage{enumitem}
\usepackage{xcolor}
\usepackage[most]{tcolorbox}
\usepackage{stmaryrd}
\usepackage[colorlinks=true,linkcolor=blue!50!black,citecolor=blue!50!black,urlcolor=blue!50!black]{hyperref}
\usepackage{fancyhdr}

\definecolor{errRed}{HTML}{B3261E}
\definecolor{warnOr}{HTML}{B26A00}
\definecolor{okGreen}{HTML}{1B6B3A}
\definecolor{infoBlue}{HTML}{1A4E8A}
\definecolor{greyed}{HTML}{555555}

\pagestyle{fancy}\fancyhf{}
\fancyhead[L]{\small Audit --- \emph{Safe prompts, safe fields} (\texttt{safe\_prompt\_fields.tex})}
\fancyhead[R]{\small\thepage}
\renewcommand{\headrulewidth}{0.4pt}

\newcommand{\RR}{\mathbb{R}}
\newcommand{\PP}{\mathbbm{P}}
\newcommand{\E}{\mathbb{E}}
\newcommand{\bth}{\boldsymbol{\theta}}
\newcommand{\be}{\mathbf{e}}
\newcommand{\bx}{\mathbf{x}}
\newcommand{\file}[1]{\texttt{#1}}
\newcommand{\lin}[1]{l.\,#1}
\newcommand{\mac}[1]{\texttt{\textbackslash #1}}
\newcommand{\by}{\mathbf{y}}
\newcommand{\cD}{\mathcal{D}}

% severity badges
\newcommand{\badge}[2]{\textcolor{#1}{\textbf{[#2]}}}
\newcommand{\Sconc}{\badge{errRed}{C: CHANGES A CONCLUSION}}
\newcommand{\Sfalse}{\badge{errRed}{F: FALSE AS STATED}}
\newcommand{\Sover}{\badge{warnOr}{O: OVERSTATED}}
\newcommand{\Srep}{\badge{warnOr}{R: REPORTING}}
\newcommand{\Sscope}{\badge{infoBlue}{S: SCOPE}}
\newcommand{\Sedit}{\badge{greyed}{E: EDITORIAL}}
\newcommand{\Ok}{\textcolor{okGreen}{\textbf{[CONFIRMED]}}}
% one-letter codes for the verdict table
\newcommand{\Tc}{\badge{errRed}{C}}\newcommand{\Tf}{\badge{errRed}{F}}\newcommand{\To}{\badge{warnOr}{O}}
\newcommand{\Tr}{\badge{warnOr}{R}}\newcommand{\Ts}{\badge{infoBlue}{S}}\newcommand{\Te}{\badge{greyed}{E}}
\newcommand{\Tok}{\textcolor{okGreen}{\textbf{[OK]}}}

\newtcolorbox{findbox}[1]{colback=white,colframe=black!25,boxrule=0.6pt,left=6pt,right=6pt,
  top=5pt,bottom=5pt,title={\small\bfseries #1},coltitle=black,colbacktitle=black!7,
  fonttitle=\bfseries,breakable}
\newtcolorbox{fixbox}[1]{colback=okGreen!4,colframe=okGreen!55,boxrule=0.6pt,left=6pt,right=6pt,
  top=5pt,bottom=5pt,title={\small\bfseries #1},coltitle=black,colbacktitle=okGreen!14,
  fonttitle=\bfseries,breakable}
\newtcolorbox{cexbox}[1]{colback=errRed!4,colframe=errRed!55,boxrule=0.6pt,left=6pt,right=6pt,
  top=5pt,bottom=5pt,title={\small\bfseries #1},coltitle=black,colbacktitle=errRed!14,
  fonttitle=\bfseries,breakable}

% numbers of the audited note (its own macro files, archived before the revision) and of the audit
\newcommand{\eoneSafePAA}{0.0024}
\newcommand{\eoneZetaA}{1.93}
\newcommand{\eonePAAa}{0.005}
\newcommand{\eonePAAb}{0.014}
\newcommand{\eonePAAc}{0.074}
\newcommand{\eonePAAd}{0.253}
\newcommand{\eoneSafePAB}{0.0046}
\newcommand{\eoneZetaB}{0.26}
\newcommand{\eonePABa}{0.007}
\newcommand{\eonePABb}{0.015}
\newcommand{\eonePABc}{0.065}
\newcommand{\eonePABd}{0.237}
\newcommand{\eoneTruePAd}{0.284}
\newcommand{\eoneOneStepPA}{0.0176}
\newcommand{\eoneCovAT}{0.897}
\newcommand{\eoneExATb}{0.07}
\newcommand{\eoneExATc}{0.24}
\newcommand{\eoneExATd}{1.00}
\newcommand{\eoneBudgetAT}{0.83}
\newcommand{\eoneCovAZ}{0.900}
\newcommand{\eoneExAZb}{0.20}
\newcommand{\eoneExAZc}{0.82}
\newcommand{\eoneExAZd}{0.95}
\newcommand{\eoneBudgetAZ}{0.50}
\newcommand{\eoneCovAR}{0.900}
\newcommand{\eoneExARb}{0.09}
\newcommand{\eoneExARc}{0.16}
\newcommand{\eoneExARd}{0.48}
\newcommand{\eoneBudgetAR}{0.98}
\newcommand{\eoneCovAE}{0.907}
\newcommand{\eoneExAEb}{0.12}
\newcommand{\eoneExAEc}{0.29}
\newcommand{\eoneExAEd}{1.00}
\newcommand{\eoneBudgetAE}{1.61}
\newcommand{\eoneCovBT}{0.900}
\newcommand{\eoneExBTb}{0.10}
\newcommand{\eoneExBTc}{0.29}
\newcommand{\eoneExBTd}{1.00}
\newcommand{\eoneBudgetBT}{0.68}
\newcommand{\eoneCovBZ}{0.901}
\newcommand{\eoneExBZb}{0.10}
\newcommand{\eoneExBZc}{0.29}
\newcommand{\eoneExBZd}{1.00}
\newcommand{\eoneBudgetBZ}{0.58}
\newcommand{\eoneCovBR}{0.900}
\newcommand{\eoneExBRb}{0.12}
\newcommand{\eoneExBRc}{0.32}
\newcommand{\eoneExBRd}{1.00}
\newcommand{\eoneBudgetBR}{0.63}
\newcommand{\eoneCovBE}{0.907}
\newcommand{\eoneExBEb}{0.12}
\newcommand{\eoneExBEc}{0.29}
\newcommand{\eoneExBEd}{1.00}
\newcommand{\eoneBudgetBE}{1.61}
\newcommand{\eoneLfourUd}{0.03}
\newcommand{\eoneLeightUd}{0.02}
\newcommand{\eoneKnnLtwoc}{0.73}
\newcommand{\eoneKnnLonec}{0.98}
\newcommand{\eoneKnnLfourc}{0.40}
\newcommand{\eoneKnnTruncc}{0.99}
\newcommand{\eoneKnnTailb}{0.29}
\newcommand{\eoneKnnTailCov}{0.918}
\newcommand{\eonePlugFail}{0.4\%}
\newcommand{\eonePlugMax}{3.7}
\newcommand{\eoneMarkovMed}{0.99}
\newcommand{\eoneRegionWorst}{0.189}
\newcommand{\eoneOneStepRectOne}{0.97}
\newcommand{\eoneMmdA}{3.3}
\newcommand{\eoneMmdB}{8.1}
\newcommand{\eoneMmdC}{34.9}
\newcommand{\eoneTwoPAasWritten}{0.142}
\newcommand{\eoneTwoPAtrunc}{0.032}
\newcommand{\eoneTwoMmdAsWritten}{84}
\newcommand{\eoneTwoMmdTrunc}{28}
\newcommand{\eoneTwoMmdNone}{149}
\newcommand{\eoneTwoBudgetUL}{0.84}
\newcommand{\eoneTwoBudgetTrunc}{0.18}
\newcommand{\eoneTwoAttackUL}{0.186}
\newcommand{\eoneTwoAttackLone}{0.054}
\newcommand{\eoneTwoAttackTrunc}{0.051}
\newcommand{\eoneTwoAttackTail}{0.044}
\newcommand{\eoneTwoAttackEU}{1.000}
\newcommand{\eoneTwoAttackEknn}{0.974}
\newcommand{\eoneTwoPlugMin}{8\%}
\newcommand{\eoneTwoPlugMax}{28\%}
\newcommand{\eoneTwoTruncMaxRatio}{1.13}
\newcommand{\eoneTwoLoneMaxRatio}{8.5}
\newcommand{\eoneTwoCertTail}{0.077}
\newcommand{\eoneTwoCertTrunc}{0.29}

% written by code/e1_walkthrough.py
\newcommand{\ewPhi}{10.1}
\newcommand{\ewRho}{1.81}
\newcommand{\ewBudRatioOne}{2.2}
\newcommand{\ewBudRatioTwo}{3.6}
\newcommand{\ewTruePAb}{0.029}
\newcommand{\ewTruePAc}{0.094}
\newcommand{\ewTruePAd}{0.283}
\newcommand{\ewMassA}{0.26}
\newcommand{\ewDispA}{1.1}
\newcommand{\ewDispNotA}{0.16}
\newcommand{\ewZetaOne}{1.81}
\newcommand{\ewOneStepAOne}{96}
\newcommand{\ewZetaTwo}{0.17}
\newcommand{\ewOneStepATwo}{3}
\newcommand{\ewSwitch}{38}
\newcommand{\ewRotA}{33}
\newcommand{\ewAngA}{-22}
\newcommand{\ewRotNear}{1.2}
\newcommand{\ewDLtwoNear}{0.08}
\newcommand{\ewDLoneNear}{0.06}
\newcommand{\ewDTrNear}{0.06}
\newcommand{\ewDTailNear}{0.00}
\newcommand{\ewDLtwoRot}{1.01}
\newcommand{\ewDLoneRot}{0.96}
\newcommand{\ewDTrRot}{0.85}
\newcommand{\ewDTailRot}{0.95}
\newcommand{\ewDLtwoUns}{0.68}
\newcommand{\ewDLoneUns}{0.49}
\newcommand{\ewDTrUns}{0.45}
\newcommand{\ewDTailUns}{0.39}
\newcommand{\ewCovU}{93}
\newcommand{\ewSafeLoU}{0.66}
\newcommand{\ewSafeHiU}{1.02}
\newcommand{\ewExUa}{5}
\newcommand{\ewExUb}{1}
\newcommand{\ewExUc}{7}
\newcommand{\ewExUd}{100}
\newcommand{\ewScUzero}{0.75}
\newcommand{\ewScUone}{0.90}
\newcommand{\ewEpsU}{0.86}
\newcommand{\ewCovK}{88}
\newcommand{\ewExKa}{15}
\newcommand{\ewExKb}{26}
\newcommand{\ewExKc}{100}
\newcommand{\ewExKd}{100}
\newcommand{\ewScKzero}{0.33}
\newcommand{\ewScKone}{1.42}
\newcommand{\ewEpsK}{0.18}
\newcommand{\ewSlicePAa}{0.256}
\newcommand{\ewSlicePAb}{0.141}
\newcommand{\ewSlicePAc}{0.028}
\newcommand{\ewSlicePAd}{0.007}
\newcommand{\ewSlicePAe}{0.002}
\newcommand{\ewSlicePAf}{0.000}
\newcommand{\ewSliceMaxU}{0.19}
\newcommand{\ewSliceMaxK}{0.034}
\newcommand{\ewKapU}{0.23}
\newcommand{\ewKapK}{0.72}
\newcommand{\ewAlphaU}{0.77}
\newcommand{\ewAlphaK}{0.28}
\newcommand{\ewRotSlice}{31}
\newcommand{\ewCertU}{0.99}
\newcommand{\ewCertRU}{1.09}
\newcommand{\ewCertQU}{0.19}
\newcommand{\ewCertPAU}{0.141}
\newcommand{\ewCertK}{0.32}
\newcommand{\ewCertRK}{0.73}
\newcommand{\ewCertQK}{0.08}
\newcommand{\ewCertPAK}{0.028}
\newcommand{\ewCertBudgetK}{0.175}
\newcommand{\ewCertFar}{0}
\newcommand{\ewAtkPrompt}{1.000}
\newcommand{\ewAtkWritten}{0.186}
\newcommand{\ewAtkImproved}{0.051}
\newcommand{\ewAtkEmbNear}{1.07}
\newcommand{\ewAtkEmbMed}{1.28}
\newcommand{\ewAtkScU}{2.3}
\newcommand{\ewAtkScK}{5.6}

% written by code/make_reflow_figure.py
\newcommand{\rfVind}{3.8}
\newcommand{\rfVrefl}{0.002}
\newcommand{\rfCostInd}{6.18}
\newcommand{\rfCostRefl}{0.97}
\newcommand{\rfCurv}{1.44}
\newcommand{\rfDecrease}{5.21}
\newcommand{\rfSum}{5.25}
\newcommand{\rfVmin}{1.6}
\newcommand{\rfVzero}{4.3}
\newcommand{\rfEoneDropped}{0.04}

% written by code/e2_figures.py
\newcommand{\etScEstatic}{0.76}
\newcommand{\etMapMaxEstatic}{0.20}
\newcommand{\etMapEstatic}{25}
\newcommand{\etMapSafeEstatic}{95}
\newcommand{\etWEtop}{0.16}
\newcommand{\etWEtopthree}{0.47}
\newcommand{\etScEadaptive}{1.37}
\newcommand{\etMapMaxEadaptive}{0.06}
\newcommand{\etMapEadaptive}{2}
\newcommand{\etMapSafeEadaptive}{100}
\newcommand{\etScFstatic}{0.77}
\newcommand{\etMapMaxTRstatic}{0.27}
\newcommand{\etMapTRstatic}{27}
\newcommand{\etMapSafeTRstatic}{89}
\newcommand{\etWFtop}{0.59}
\newcommand{\etWFtopthree}{0.87}
\newcommand{\etScFadaptive}{1.57}
\newcommand{\etMapMaxTRadaptive}{0.04}
\newcommand{\etMapTRadaptive}{0}
\newcommand{\etMapSafeTRadaptive}{95}
\newcommand{\etTuneGap}{0.05}
\newcommand{\etTuneRegret}{0.005}
\newcommand{\etNearGap}{0.004}
\newcommand{\etDensEA}{0.16}
\newcommand{\etDensUEA}{0.12}
\newcommand{\etEpsEA}{0.63}
\newcommand{\etEpsUEA}{1.60}
\newcommand{\etDensEB}{0.99}
\newcommand{\etDensUEB}{0.13}
\newcommand{\etEpsEB}{0.22}
\newcommand{\etEpsUEB}{1.57}
\newcommand{\etDensOneE}{0.17}
\newcommand{\etDensFA}{0.20}
\newcommand{\etDensUFA}{0.13}
\newcommand{\etEpsFA}{0.20}
\newcommand{\etEpsUFA}{0.59}
\newcommand{\etDensFB}{1.00}
\newcommand{\etDensUFB}{0.15}
\newcommand{\etEpsFB}{0.08}
\newcommand{\etEpsUFB}{0.58}
\newcommand{\etDensOneF}{0.24}
\newcommand{\etDensPA}{0.21}
\newcommand{\etDensUPA}{0.14}
\newcommand{\etEpsPA}{0.15}
\newcommand{\etEpsUPA}{0.58}
\newcommand{\etDensPB}{1.00}
\newcommand{\etDensUPB}{0.16}
\newcommand{\etEpsPB}{0.07}
\newcommand{\etEpsUPB}{0.57}
\newcommand{\etDensOneP}{0.26}
\newcommand{\etAtkOneES}{1.000}
\newcommand{\etAtkMedOneES}{0.983}
\newcommand{\etEnfOneES}{0.049}
\newcommand{\etMmdOneES}{28}
\newcommand{\etAtkOneEA}{0.976}
\newcommand{\etAtkMedOneEA}{0.920}
\newcommand{\etEnfOneEA}{0.009}
\newcommand{\etMmdOneEA}{11}
\newcommand{\etAtkOnePBS}{0.262}
\newcommand{\etAtkMedOnePBS}{0.131}
\newcommand{\etEnfOnePBS}{0.132}
\newcommand{\etMmdOnePBS}{73}
\newcommand{\etAtkOnePBA}{0.030}
\newcommand{\etAtkMedOnePBA}{0.021}
\newcommand{\etEnfOnePBA}{0.019}
\newcommand{\etMmdOnePBA}{21}
\newcommand{\etAtkOneFS}{0.220}
\newcommand{\etAtkMedOneFS}{0.146}
\newcommand{\etEnfOneFS}{0.077}
\newcommand{\etMmdOneFS}{49}
\newcommand{\etAtkOneFA}{0.036}
\newcommand{\etAtkMedOneFA}{0.020}
\newcommand{\etEnfOneFA}{0.026}
\newcommand{\etMmdOneFA}{26}
\newcommand{\etAtkTwoES}{0.655}
\newcommand{\etAtkMedTwoES}{0.537}
\newcommand{\etEnfTwoES}{0.061}
\newcommand{\etMmdTwoES}{78}
\newcommand{\etAtkTwoEA}{0.496}
\newcommand{\etAtkMedTwoEA}{0.228}
\newcommand{\etEnfTwoEA}{0.027}
\newcommand{\etMmdTwoEA}{47}
\newcommand{\etAtkTwoPBS}{0.099}
\newcommand{\etAtkMedTwoPBS}{0.086}
\newcommand{\etEnfTwoPBS}{0.057}
\newcommand{\etMmdTwoPBS}{89}
\newcommand{\etAtkTwoPBA}{0.030}
\newcommand{\etAtkMedTwoPBA}{0.027}
\newcommand{\etEnfTwoPBA}{0.034}
\newcommand{\etMmdTwoPBA}{76}
\newcommand{\etAtkTwoFS}{0.072}
\newcommand{\etAtkMedTwoFS}{0.064}
\newcommand{\etEnfTwoFS}{0.040}
\newcommand{\etMmdTwoFS}{75}
\newcommand{\etAtkTwoFA}{0.039}
\newcommand{\etAtkMedTwoFA}{0.029}
\newcommand{\etEnfTwoFA}{0.034}
\newcommand{\etMmdTwoFA}{72}
\newcommand{\etExESa}{0.11}
\newcommand{\etExESb}{0.12}
\newcommand{\etExESc}{0.30}
\newcommand{\etExESd}{1.00}
\newcommand{\etJES}{0.38}
\newcommand{\etCovES}{0.904}
\newcommand{\etExEAa}{0.11}
\newcommand{\etExEAb}{0.16}
\newcommand{\etExEAc}{0.99}
\newcommand{\etExEAd}{1.00}
\newcommand{\etJEA}{0.57}
\newcommand{\etCovEA}{0.905}
\newcommand{\etBestE}{0.3}
\newcommand{\etExFSa}{0.12}
\newcommand{\etExFSb}{0.13}
\newcommand{\etExFSc}{0.35}
\newcommand{\etExFSd}{0.98}
\newcommand{\etJFS}{0.40}
\newcommand{\etCovFS}{0.904}
\newcommand{\etExFAa}{0.11}
\newcommand{\etExFAb}{0.20}
\newcommand{\etExFAc}{0.99}
\newcommand{\etExFAd}{1.00}
\newcommand{\etJFA}{0.58}
\newcommand{\etCovFA}{0.906}
\newcommand{\etBestF}{0.2}
\newcommand{\etExPSa}{0.13}
\newcommand{\etExPSb}{0.14}
\newcommand{\etExPSc}{0.32}
\newcommand{\etExPSd}{1.00}
\newcommand{\etJPS}{0.40}
\newcommand{\etCovPS}{0.903}
\newcommand{\etExPAa}{0.11}
\newcommand{\etExPAb}{0.21}
\newcommand{\etExPAc}{1.00}
\newcommand{\etExPAd}{1.00}
\newcommand{\etJPA}{0.58}
\newcommand{\etCovPA}{0.907}
\newcommand{\etBestP}{0.1}
\newcommand{\etExSSa}{0.08}
\newcommand{\etExSSb}{0.15}
\newcommand{\etExSSc}{0.83}
\newcommand{\etExSSd}{0.91}
\newcommand{\etJSS}{0.49}
\newcommand{\etCovSS}{0.902}
\newcommand{\etExSAa}{0.13}
\newcommand{\etExSAb}{0.59}
\newcommand{\etExSAc}{1.00}
\newcommand{\etExSAd}{1.00}
\newcommand{\etJSA}{0.68}
\newcommand{\etCovSA}{0.905}
\newcommand{\etBestS}{0.2}
\newcommand{\etExLSa}{0.08}
\newcommand{\etExLSb}{0.08}
\newcommand{\etExLSc}{0.25}
\newcommand{\etExLSd}{0.99}
\newcommand{\etJLS}{0.35}
\newcommand{\etCovLS}{0.902}
\newcommand{\etExLAa}{0.09}
\newcommand{\etExLAb}{0.13}
\newcommand{\etExLAc}{0.74}
\newcommand{\etExLAd}{1.00}
\newcommand{\etJLA}{0.49}
\newcommand{\etCovLA}{0.902}
\newcommand{\etBestL}{0.1}
\newcommand{\etExTSa}{0.12}
\newcommand{\etExTSb}{0.14}
\newcommand{\etExTSc}{0.53}
\newcommand{\etExTSd}{1.00}
\newcommand{\etJTS}{0.45}
\newcommand{\etCovTS}{0.905}
\newcommand{\etExTAa}{0.12}
\newcommand{\etExTAb}{0.24}
\newcommand{\etExTAc}{0.99}
\newcommand{\etExTAd}{1.00}
\newcommand{\etJTA}{0.59}
\newcommand{\etCovTA}{0.904}
\newcommand{\etBestT}{0.3}
\newcommand{\etSpreadOne}{0.48}
\newcommand{\etDispOne}{0.24}
\newcommand{\etSpreadTwo}{0.39}
\newcommand{\etDispTwo}{0.21}
\newcommand{\etCovMin}{0.897}
\newcommand{\etCovMax}{0.908}
\newcommand{\etPoolPAOne}{0.255}
\newcommand{\etPoolPATwo}{0.208}
\newcommand{\etNuCovMin}{0.895}
\newcommand{\etNuCovMax}{0.908}
\newcommand{\etNuSafePA}{0.0024}
\newcommand{\etNuPAd}{0.255}
\newcommand{\etNuMmd}{3.4}
\newcommand{\etNuSame}{0.054}
\newcommand{\etNuPairs}{0.471}
\newcommand{\etNuEzeroS}{0.30}
\newcommand{\etNuEzeroA}{0.99}
\newcommand{\etNuEtwoS}{0.21}
\newcommand{\etNuEtwoA}{0.18}
\newcommand{\etNuEfourS}{0.15}
\newcommand{\etNuEfourA}{0.13}
\newcommand{\etNuFS}{0.31}
\newcommand{\etNuFA}{0.97}
\newcommand{\etNuPS}{0.32}
\newcommand{\etNuPA}{0.93}

\newcommand{\pairKLAmed}{0.067}
\newcommand{\pairKLAninety}{0.257}
\newcommand{\pALmean}{8.3\times 10^{-6}}
\newcommand{\pairKLDmed}{0.378}
\newcommand{\pairKLDninety}{0.755}
\newcommand{\pALDmean}{0.0004}
\newcommand{\pALCmean}{0.0017}
\newcommand{\covAF}{0.897}
\newcommand{\covAK}{0.898}
\newcommand{\exclAFthree}{0.42}
\newcommand{\exclAKthree}{1.00}
\newcommand{\exclAKone}{0.83}
\newcommand{\exclAFone}{0.19}
\newcommand{\exclCrefthree}{0.15}
\newcommand{\exclConthree}{0.55}
\newcommand{\exclDfwdthree}{0.17}
\newcommand{\exclDrevthree}{0.55}
\newcommand{\exclDfwdten}{0.69}
\newcommand{\certAKmed}{0.031}
\newcommand{\certAFmed}{0.283}
\newcommand{\certDrevmed}{0.134}
\newcommand{\certDfwdmed}{0.553}
\newcommand{\insideAKmed}{0.0123}
\newcommand{\insideAFmed}{0.0324}
\newcommand{\insideAFmax}{0.0947}
\newcommand{\insideDrevmed}{0.0145}
\newcommand{\insideDfwdmed}{0.0707}
\newcommand{\geoCrefclean}{0.28}
\newcommand{\geoConclean}{0.35}
\newcommand{\winCrefclean}{32}
\newcommand{\winConclean}{32}
\newcommand{\geoCrefcont}{0.21}
\newcommand{\geoConcont}{0.20}
\newcommand{\winCrefcont}{32}
\newcommand{\winConcont}{32}
\newcommand{\ntasks}{32}
\newcommand{\pAUcont}{0.138}
\newcommand{\pACrefcont}{0.047}
\newcommand{\pAConcont}{0.025}
\newcommand{\pASten}{0.0087}
\newcommand{\backoffmax}{0.011}
\newcommand{\backofffrac}{0.70}
\newcommand{\socptime}{0.31}
\newcommand{\epsref}{0.0182}
\newcommand{\epson}{0.0193}
\newcommand{\exampletask}{8}
\newcommand{\exUpA}{0.133}
\newcommand{\exCrefpA}{0.029}
\newcommand{\exConpA}{0.013}
\newcommand{\exSpA}{0.009}
\newcommand{\excentrepA}{0.001}
\newcommand{\exoraclepA}{0.003}
\newcommand{\excentresref}{0.55}
\newcommand{\excentreson}{0.54}
\newcommand{\lipmin}{52}
\newcommand{\lipmax}{53}
\newcommand{\gronwalllog}{15}

% written by code/audit_checks.py
\newcommand{\auMaxProbeOne}{16.5}
\newcommand{\auMaxProbeTwo}{7.6}
\newcommand{\auMaxEvalOne}{2.8\times10^{4}}
\newcommand{\auMaxEvalTwo}{6.1}
\newcommand{\auMaxGridOne}{2.1\times10^{7}}
\newcommand{\auMaxReflow}{2.5\times10^{5}}
\newcommand{\auProbe}{136}
\newcommand{\auProbeNorm}{3.96}
\newcommand{\auBlowSafe}{44}
\newcommand{\auBlowa}{55}
\newcommand{\auBlowb}{57}
\newcommand{\auBlowc}{34}
\newcommand{\auBlowd}{3}
\newcommand{\auShareTwoMed}{6}
\newcommand{\auShareTwoMean}{19}
\newcommand{\auShareTwoNinety}{57}
\newcommand{\auShareFourMed}{35}
\newcommand{\auShareFourMean}{45}
\newcommand{\auShareFourNinety}{99}
\newcommand{\auOneLfourUAlla}{0.17}
\newcommand{\auOneLfourUAllb}{0.17}
\newcommand{\auOneLfourUAllc}{0.08}
\newcommand{\auOneLfourUAlld}{0.03}
\newcommand{\auOneLfourUAllEps}{1.816}
\newcommand{\auOneLfourUDropa}{0.07}
\newcommand{\auOneLfourUDropb}{0.07}
\newcommand{\auOneLfourUDropc}{0.30}
\newcommand{\auOneLfourUDropd}{1.00}
\newcommand{\auOneLfourUDropEps}{0.911}
\newcommand{\auOneLeightUAlla}{0.18}
\newcommand{\auOneLeightUAllb}{0.17}
\newcommand{\auOneLeightUAllc}{0.08}
\newcommand{\auOneLeightUAlld}{0.02}
\newcommand{\auOneLeightUAllEps}{3.522}
\newcommand{\auOneLeightUDropa}{0.08}
\newcommand{\auOneLeightUDropb}{0.08}
\newcommand{\auOneLeightUDropc}{0.18}
\newcommand{\auOneLeightUDropd}{1.00}
\newcommand{\auOneLeightUDropEps}{1.163}
\newcommand{\auOneLtwoNAlla}{0.13}
\newcommand{\auOneLtwoNAllb}{0.17}
\newcommand{\auOneLtwoNAllc}{0.73}
\newcommand{\auOneLtwoNAlld}{1.00}
\newcommand{\auOneLtwoNAllEps}{0.377}
\newcommand{\auOneLtwoNDropa}{0.11}
\newcommand{\auOneLtwoNDropb}{0.20}
\newcommand{\auOneLtwoNDropc}{1.00}
\newcommand{\auOneLtwoNDropd}{1.00}
\newcommand{\auOneLtwoNDropEps}{0.239}
\newcommand{\auOneLtwoUAlla}{0.09}
\newcommand{\auOneLtwoUAllb}{0.07}
\newcommand{\auOneLtwoUAllc}{0.24}
\newcommand{\auOneLtwoUAlld}{1.00}
\newcommand{\auOneLtwoUAllEps}{0.829}
\newcommand{\auOneLtwoUDropa}{0.07}
\newcommand{\auOneLtwoUDropb}{0.08}
\newcommand{\auOneLtwoUDropc}{0.33}
\newcommand{\auOneLtwoUDropd}{1.00}
\newcommand{\auOneLtwoUDropEps}{0.739}
\newcommand{\auOneLfourNAlla}{0.14}
\newcommand{\auOneLfourNAllb}{0.14}
\newcommand{\auOneLfourNAllc}{0.40}
\newcommand{\auOneLfourNAlld}{0.89}
\newcommand{\auOneLfourNAllEps}{0.660}
\newcommand{\auOneLfourNDropa}{0.09}
\newcommand{\auOneLfourNDropb}{0.15}
\newcommand{\auOneLfourNDropc}{1.00}
\newcommand{\auOneLfourNDropd}{1.00}
\newcommand{\auOneLfourNDropEps}{0.389}
\newcommand{\auTwoLfourUAlla}{0.10}
\newcommand{\auTwoLfourUAllb}{0.10}
\newcommand{\auTwoLfourUAllc}{0.17}
\newcommand{\auTwoLfourUAlld}{0.97}
\newcommand{\auTwoLfourUAllEps}{0.967}
\newcommand{\auTwoLfourUDropa}{0.08}
\newcommand{\auTwoLfourUDropb}{0.09}
\newcommand{\auTwoLfourUDropc}{0.30}
\newcommand{\auTwoLfourUDropd}{1.00}
\newcommand{\auTwoLfourUDropEps}{0.769}
\newcommand{\auTwoLeightUAlla}{0.10}
\newcommand{\auTwoLeightUAllb}{0.08}
\newcommand{\auTwoLeightUAllc}{0.07}
\newcommand{\auTwoLeightUAlld}{0.23}
\newcommand{\auTwoLeightUAllEps}{1.715}
\newcommand{\auTwoLeightUDropa}{0.08}
\newcommand{\auTwoLeightUDropb}{0.09}
\newcommand{\auTwoLeightUDropc}{0.25}
\newcommand{\auTwoLeightUDropd}{1.00}
\newcommand{\auTwoLeightUDropEps}{0.954}
\newcommand{\auTwoLtwoNAlla}{0.10}
\newcommand{\auTwoLtwoNAllb}{0.19}
\newcommand{\auTwoLtwoNAllc}{1.00}
\newcommand{\auTwoLtwoNAlld}{1.00}
\newcommand{\auTwoLtwoNAllEps}{0.188}
\newcommand{\auTwoLtwoNDropa}{0.11}
\newcommand{\auTwoLtwoNDropb}{0.19}
\newcommand{\auTwoLtwoNDropc}{1.00}
\newcommand{\auTwoLtwoNDropd}{1.00}
\newcommand{\auTwoLtwoNDropEps}{0.178}
\newcommand{\auTwoLtwoUAlla}{0.08}
\newcommand{\auTwoLtwoUAllb}{0.10}
\newcommand{\auTwoLtwoUAllc}{0.29}
\newcommand{\auTwoLtwoUAlld}{1.00}
\newcommand{\auTwoLtwoUAllEps}{0.676}
\newcommand{\auTwoLtwoUDropa}{0.07}
\newcommand{\auTwoLtwoUDropb}{0.09}
\newcommand{\auTwoLtwoUDropc}{0.30}
\newcommand{\auTwoLtwoUDropd}{1.00}
\newcommand{\auTwoLtwoUDropEps}{0.650}
\newcommand{\auTwoLfourNAlla}{0.10}
\newcommand{\auTwoLfourNAllb}{0.13}
\newcommand{\auTwoLfourNAllc}{0.97}
\newcommand{\auTwoLfourNAlld}{1.00}
\newcommand{\auTwoLfourNAllEps}{0.286}
\newcommand{\auTwoLfourNDropa}{0.11}
\newcommand{\auTwoLfourNDropb}{0.21}
\newcommand{\auTwoLfourNDropc}{1.00}
\newcommand{\auTwoLfourNDropd}{1.00}
\newcommand{\auTwoLfourNDropEps}{0.227}
\newcommand{\auELtwoNAlla}{0.09}
\newcommand{\auELtwoNAllb}{0.13}
\newcommand{\auELtwoNAllc}{0.74}
\newcommand{\auELtwoNAlld}{1.00}
\newcommand{\auELtwoNDropa}{0.11}
\newcommand{\auELtwoNDropb}{0.22}
\newcommand{\auELtwoNDropc}{1.00}
\newcommand{\auELtwoNDropd}{1.00}
\newcommand{\auETrNAlla}{0.11}
\newcommand{\auETrNAllb}{0.20}
\newcommand{\auETrNAllc}{0.99}
\newcommand{\auETrNAlld}{1.00}
\newcommand{\auELtwoUAlla}{0.08}
\newcommand{\auELtwoUAllb}{0.08}
\newcommand{\auELtwoUAllc}{0.25}
\newcommand{\auELtwoUAlld}{0.99}
\newcommand{\auELtwoUDropa}{0.13}
\newcommand{\auELtwoUDropb}{0.14}
\newcommand{\auELtwoUDropc}{0.39}
\newcommand{\auELtwoUDropd}{1.00}
\newcommand{\auETrUAlla}{0.12}
\newcommand{\auETrUAllb}{0.13}
\newcommand{\auETrUAllc}{0.35}
\newcommand{\auETrUAlld}{0.98}
\newcommand{\auETrUDropa}{0.12}
\newcommand{\auETrUDropb}{0.13}
\newcommand{\auETrUDropc}{0.35}
\newcommand{\auETrUDropd}{0.98}
\newcommand{\auETrNDropa}{0.11}
\newcommand{\auETrNDropb}{0.20}
\newcommand{\auETrNDropc}{0.99}
\newcommand{\auETrNDropd}{1.00}
\newcommand{\auProbeCovMin}{0.897}
\newcommand{\auProbeCovMax}{0.934}
\newcommand{\auQminOne}{1.5\times10^{-3}}
\newcommand{\auQmaxOne}{2.4\times10^{4}}
\newcommand{\auQcondOne}{1.6\times10^{7}}
\newcommand{\auQminTwo}{4.3\times10^{-4}}
\newcommand{\auQmaxTwo}{2.9\times10^{2}}
\newcommand{\auQcondTwo}{6.7\times10^{5}}
\newcommand{\auPBoneAtkOneS}{0.262}
\newcommand{\auPBoneAtkTwoS}{0.85}
\newcommand{\auTRoneAtkTwoS}{0.071}
\newcommand{\auPBtwoAtkTwoS}{0.049}
\newcommand{\auTRtwoAtkTwoS}{0.047}
\newcommand{\auMHAtkOneS}{0.003}
\newcommand{\auMHAtkTwoS}{0.026}
\newcommand{\auMHCovS}{0.905}
\newcommand{\auMHRhoS}{1\times10^{-6}}
\newcommand{\auMHExSa}{1.00}
\newcommand{\auMHExSb}{1.00}
\newcommand{\auMHExSc}{1.00}
\newcommand{\auMHExSd}{1.00}
\newcommand{\auPBoneAtkOneA}{0.030}
\newcommand{\auPBoneAtkTwoA}{0.22}
\newcommand{\auTRoneAtkTwoA}{0.019}
\newcommand{\auPBtwoAtkTwoA}{0.024}
\newcommand{\auTRtwoAtkTwoA}{0.029}
\newcommand{\auMHAtkOneA}{0.004}
\newcommand{\auMHAtkTwoA}{0.044}
\newcommand{\auMHCovA}{0.905}
\newcommand{\auMHRhoA}{1\times10^{-6}}
\newcommand{\auMHExAa}{1.00}
\newcommand{\auMHExAb}{1.00}
\newcommand{\auMHExAc}{1.00}
\newcommand{\auMHExAd}{1.00}
\newcommand{\auMHcAd}{0.03}
\newcommand{\auMHEnfS}{0.001}
\newcommand{\auMHMmdS}{91}
\newcommand{\auMHEnfA}{0.001}
\newcommand{\auMHMmdA}{101}
\newcommand{\auFieldAtkMin}{0.036}
\newcommand{\auFieldAtkMax}{0.220}
\newcommand{\auEoneAtkTwoMax}{0.084}
\newcommand{\auEoneAtkTwoNearMax}{0.024}
\newcommand{\auMvar}{0}
\newcommand{\auMHNuMin}{1.00}
\newcommand{\auBudgetOne}{0.83}
\newcommand{\auFloorOne}{0.33}
\newcommand{\auTypicalOne}{0.68}
\newcommand{\auBudgetTwo}{0.68}
\newcommand{\auFloorTwo}{0.27}
\newcommand{\auTypicalTwo}{0.54}
\newcommand{\auFigCovMin}{0.897}
\newcommand{\auFigCovMax}{0.934}
\newcommand{\auTailCovOne}{0.918}
\newcommand{\auTailCovTwo}{0.934}
\newcommand{\auTailProbesOne}{11}
\newcommand{\auTailProbesTwo}{2}
\newcommand{\auTailOneCovTwo}{0.958}
\newcommand{\auTailOneEpsTwo}{0.000}
\newcommand{\auTruncCovOne}{0.908}
\newcommand{\auSlackTrBound}{0.300}
\newcommand{\auSlackTrPA}{0.030}
\newcommand{\auBudgetTr}{0.182}
\newcommand{\auSlackTailBound}{0.086}
\newcommand{\auSlackTailPA}{0.026}
\newcommand{\auBudgetTail}{0.047}
\newcommand{\auHoeff}{0.027}
\newcommand{\auHoeffTailPct}{58}
\newcommand{\auHoeffTrPct}{15}
\newcommand{\auCPzero}{0.0015}
\newcommand{\auRegionBoundTr}{0.34}
\newcommand{\auRegionBoundTail}{0.09}
\newcommand{\auOneStepPA}{0.0176}
\newcommand{\auRectOnePA}{0.0024}
\newcommand{\auOneStepRatio}{7}
\newcommand{\auExcessa}{0.0023}
\newcommand{\auTargeta}{0.0095}
\newcommand{\auRealiseda}{24}
\newcommand{\auExcessb}{0.0112}
\newcommand{\auTargetb}{0.0284}
\newcommand{\auRealisedb}{40}
\newcommand{\auLipMin}{52}
\newcommand{\auGronLog}{15}
\newcommand{\auExpLLog}{22.6}
\newcommand{\auPocOnOne}{0.133}
\newcommand{\auPocOnThree}{0.552}
\newcommand{\auPocRevOne}{0.129}
\newcommand{\auPocRevThree}{0.546}
\newcommand{\auMixcheck}{5.7\times10^{-13}}
\newcommand{\auSpreadOld}{0.485}
\newcommand{\auSpreadNew}{0.484}
\newcommand{\auDispOld}{0.236}
\newcommand{\auDispNew}{0.238}
\newcommand{\auVerOverOneTr}{0.0\%}
\newcommand{\auVerN}{5600}
\newcommand{\auVerViol}{0}

\makeatletter\newcommand{\inputrows}[1]{\@@input #1 }\makeatother

\title{\vspace{-1.2cm}\textbf{Audit of \emph{Safe prompts, safe fields}}\\[2pt]
\large the order-$1$ Wasserstein conformal region for linear flow matching\\[6pt]
\normalsize \file{flow\_matching\_safe\_set\_POC/safe\_prompt\_fields.tex}, version of 14 October 2026 (29 pages)}
\author{\normalsize Mathematical and experimental audit}
\date{\normalsize 14 October 2026}

\begin{document}
\maketitle
\thispagestyle{fancy}
\vspace{-0.6cm}

\begin{abstract}\noindent
Every numbered statement of the note (Sections~2--7 and Appendix~A) was re-derived line by line, every script of
the two experiments of Sections~9 and~10 (\file{code/e1\_*.py}, \file{code/e2\_*.py}, \file{code/make\_reflow\_figure.py})
was read against the text, and the numbers that carry a conclusion were recomputed from \file{results/}. Three new
checks were run on the same pools and models, without refitting anything: a probe-sensitivity analysis, a second
attack that searches the directions a region penalises least, and a Mahalanobis prompt-space baseline
(\file{code/e1\_robustness.py}, \file{code/e2\_robustness.py}; every number below is produced by
\file{code/audit\_checks.py}).
\smallskip

\noindent\textbf{Headline.} The theory is sound: every theorem, proposition and lemma is correct as stated, and its
proof is complete. The defects are in the experiments and in the prose that draws conclusions from them.
Three experimental conclusions do not survive re-checking. (F1) The finding that $L^2$ and $L^p$ map distances are less
sensitive than the truncated $L^1$ distance is produced by \emph{one} of the 256 probe noises, which the 1-rectified
sampler sends far from the data; without it, $L^2$ is at least as sensitive, as Proposition~7.3 predicts. (F2) The
pull-back region of the 1-rectified model is broken by an attack that searches the low-eigenvalue directions of its
metric ($P(A)=\auPBoneAtkTwoA$ instead of the reported $\etAtkOnePBA$). (F3) A Mahalanobis prompt-space region excludes
every unsafe request and resists the note's attack as well as any field region, because the approved embeddings of the
toy are degenerate; the note's ``prompt versus field'' conclusion is shown only against Euclidean prompt distances.
\smallskip

\noindent\textbf{Other defects.} Two statements in the prose are false (convexity of the improved region, a ``Markov-free'' certificate) and
three overstated (the size of the Gr\"onwall factor, the mechanism attributed to the floor of Proposition~5.6, a
comparison of sensitivities), ten reporting or methodology defects (among them a misreported coverage range, a certificate check
that cannot fail, and an untested verification step presented as part of the method), and seven editorial items.
What survives: coverage in every experiment, the failure of static weights on a diverse family, the repair by
nearest-anchor weights (confirmed on disjoint tuning data), and the robustness of field-space regions to both attacks.
\end{abstract}

\vspace{-0.2cm}
\section{Scope, conventions and method}

\paragraph{Scope.} The note \file{safe\_prompt\_fields.tex} as of 14 October 2026, 10:10 (1954 lines, 29 pages),
archived unchanged with its number macros, its README and the three figure scripts that the revision modifies in
\file{archive/2026-10-14\_pre\_audit/}. Line numbers (\lin{n}) refer to that archived file; result, table and figure
numbers to its compiled PDF. Experiments: E1 = Section~9 (pool of \file{e1\_fit.py}, 300 safe and 400 unsafe prompts),
E2 = Section~10 (pool of \file{e2\_pool.py}). ``1-rectified'' and ``reflowed'' are the two models of E1; $P(A)$ is the
fraction of outputs in the unsafe disk $A=\{\|\by\|\le0.6\}$.

\paragraph{Severity scale.}
\begin{itemize}[leftmargin=1.4em,itemsep=1pt,topsep=2pt]
\item \Sconc{} --- a conclusion of the note does not hold once the defect is removed.
\item \Sfalse{} --- the printed statement admits a counterexample.
\item \Sover{} --- the claim goes beyond its evidence or rests on a wrong mechanism, but its direction may be right.
\item \Srep{} --- a reported number, check or protocol does not show what the text says it shows.
\item \Sscope{} --- a limitation that the note does not state.
\item \Sedit{} --- wording or bookkeeping; no change of content.
\item \Ok{} --- statement and proof (or number and its script) verified, no action needed.
\end{itemize}

\paragraph{Method.}
\emph{(i) Mathematics}: every proof of Sections~2--7 and Appendix~A re-derived, including the measure-theoretic
details (versions of conditional expectations, null sets at $t=0$, well-posedness), and the cited theorem numbers
checked against the bibliography's notes. \emph{(ii) Code}: every script of E1 and E2 read against the protocol stated
in the text, including the random-number streams, the noise sets used on each side of each inequality, and the pools
used for tuning and evaluation. \emph{(iii) Numbers}: the quoted numbers re-derived from \file{results/}; where a
conclusion looked fragile, a targeted check was designed and run (Section~\ref{sec:repro}).

\clearpage
\section{Verdict table}

Codes as in the severity scale: \Tc{} changes a conclusion, \Tf{} false as stated, \To{} overstated, \Tr{} reporting,
\Ts{} scope, \Te{} editorial, \Tok{} confirmed.
\smallskip

\renewcommand{\arraystretch}{1.18}
\begin{longtable}{@{}p{0.9cm}>{\raggedright\arraybackslash}p{3.3cm}p{1.0cm}>{\raggedright\arraybackslash}p{10.3cm}@{}}
\toprule
\textbf{Id} & \textbf{Location} & \textbf{Code} & \textbf{Summary}\\
\midrule
\endhead
\multicolumn{4}{@{}l}{\textbf{Findings that change a conclusion}}\\
\midrule
F1 & \S9 Steps 4--5, Fig.~8, \S10.1 & \Tc & One probe noise (of 256), sent beyond radius $4$ by the 1-rectified sampler, produces ``$L^p$, $p>2$, fails'' and ``$L^2$ is less sensitive than truncated $L^1$''.\\
F2 & \S10.4, Table~8 & \Tc & The pull-back region of the 1-rectified model admits prompts with $P(A)=\auPBoneAtkTwoA$ (adaptive) and $\auPBoneAtkTwoS$ (static); reported: $\etAtkOnePBA$, $\etAtkOnePBS$.\\
F3 & abstract, \S9 Step 8, \S10.3, Q2 & \Tc & A Mahalanobis prompt region excludes every unsafe request and resists the attack; ``fields, not prompts'' is shown only against Euclidean prompt distances.\\
\midrule
\multicolumn{4}{@{}l}{\textbf{Mathematics}}\\
\midrule
M1 & \lin{1553}, \lin{1880--1881} & \Tf & ``One-step models make [the improved region] convex'': a union of balls, and $\min(\|\cdot\|,r)$ is not convex.\\
M2 & abstract \lin{108}, \lin{1084--1085}, \lin{1552}, \lin{1871} & \Tf & ``Needs no Markov step'': the proof of Proposition~7.5 is Markov's inequality applied to $\min(\delta,r)$.\\
M3 & \lin{149--150}, \lin{379--380} & \To & ``$e^L$ is about $10^{15}$'': $10^{15}$ is a discrete product with an a-priori bound $L\approx\auLipMin$; $e^{\auLipMin}\approx10^{\auExpLLog}$.\\
M4 & \lin{1350--1353}, \lin{1686--1687}, \lin{1807} & \To & Detection failure attributed to the floor of Proposition~5.6(iii); the budget is $2.5$ times the floor.\\
M5 & \lin{1043--1048} & \To & ``The $L^2$ map metric is the most sensitive'' compares lower bounds across distances whose thresholds differ.\\
M6 & \lin{645}, \lin{655--656}, \lin{1389--1394}, \lin{1551} & \Te & ``Fidelity can only be lost''; Lemma~3.9(iii) applied to trimmed pairs; affine combinations called mixtures; ``order-$1$ structure'' kept in name only.\\
\midrule
\multicolumn{4}{@{}l}{\textbf{Methodology and reporting}}\\
\midrule
R1 & \lin{1359}, Table~4, Fig.~8 & \Tr & ``Coverage between $\eoneCovAT$ and $\eoneCovAE$ for every distance'': the tail score has $\auTailCovOne$ and $\auTailCovTwo$ (ties).\\
R2 & Table~5, \lin{1438} & \Tr & ``No enforced field violates its certificate'': both sides are measured on the same noise; the check cannot fail at this slack.\\
R3 & \lin{1440--1445}, \lin{1545--1551} & \Tr & The fresh-noise verification is part of ``the improved method'' but was never run; probability qualifier missing; Hoeffding margin eats $\auHoeffTailPct\%$ of the tail budget.\\
R4 & abstract \lin{108--109} & \To & ``Field-space regions cap the unsafe mass at a few percent'': the region as written admits $\eoneTwoAttackUL$ (E1) and $\etAtkOneFS$ (E2).\\
R5 & \lin{1859--1860} & \Tr & Q3 omits that the reflowed one-step model puts $\auOneStepPA$ of approved outputs in $A$ ($\auOneStepRatio\times$ the 1-rectified model).\\
R6 & \lin{1500--1501}, \lin{1671--1672}, \lin{1817--1818} & \To & Small-mass detection limits are confounded: the model realises only $\auRealiseda\%$ ($m=0.01$) and $\auRealisedb\%$ ($m=0.03$) of the requested unsafe mass.\\
R7 & \S9, \S10 & \Tr & No uncertainty is reported; attacks on one split; E1's variants selected on the pool they are reported on.\\
R8 & \S9--\S11 & \Ts & Nearest-anchor detection depends on the safe family being two-dimensional; not discussed.\\
R9 & \S7, \S9, title & \Ts & The improved method is a nearest-neighbour conformal novelty detector; related work absent.\\
R10 & \S11 \lin{1873} & \Tr & ``Tested: prompt mixing'': only the algebraic identity was checked; the projection's results were computed and not reported.\\
\midrule
\multicolumn{4}{@{}l}{\textbf{Editorial and bookkeeping}}\\
\midrule
H1--H7 & README, \file{e2\_figures.py}, Table~2, \lin{11} & \Te & Stale page count, two numbers computed on the wrong pools, an unexplained coincidence in Table~2, a stale header comment, unused macros, a non-i.i.d.\ Clopper--Pearson bound.\\
\midrule
\multicolumn{4}{@{}l}{\textbf{Confirmed}}\\
\midrule
OK & \S2--\S7, App.~A, all numbered results & \Tok & Statements and proofs correct; details in Section~\ref{sec:ok}.\\
\bottomrule
\end{longtable}

% =====================================================================================
\section{Findings that change a conclusion}\label{sec:major}
% =====================================================================================

\begin{findbox}{F1 \quad One probe noise drives the comparison of map distances \hfill \Sconc}
\textbf{What the note says.} Section~9, Step~4 (\lin{1322--1327}): ``the $L^p$ distances with $p>2$ [\dots] fail: high powers
weigh the largest displacements, and between safe prompts the largest displacements are blob changes, which exceed the
displacement of an unsafe mode. With uniform weights they exclude only $\eoneLfourUd$ ($p=4$) and $\eoneLeightUd$ ($p=8$) of the
prompts with $m=0.3$.'' Step~5 (\lin{1359--1362}): exclusion at $m=0.1$ rises from $\eoneExATc$ to $\eoneKnnLtwoc$ with the nearest field in
$L^2$, but to $\eoneKnnLonec$ in $L^1$ and $\eoneKnnTruncc$ in truncated $L^1$. Figure~8 caption (\lin{1382--1383}): the nearest field works ``with
$L^2$ (orange) only on the reflowed model''. Section~10.1 (\lin{1591--1592}): the field region uses ``the truncated map
distance, the most sensitive field distance of Section~9''.

\smallskip\textbf{What the data show.}
\begin{itemize}[leftmargin=1.3em,itemsep=1pt,topsep=2pt]
\item The 50-step Euler sampler of the 1-rectified model sends probe~$\auProbe$ ($|\bx_0|=\auProbeNorm$, the largest of the
256 frozen probes) beyond radius~$4$ for $\auBlowSafe\%$ of the safe prompts, up to $|T|=\auMaxProbeOne$; the ring has radius
about~$2$. No other probe leaves radius~$4$ for any pool prompt, in E1 or in E2, and the same probe is the only one flagged for
the reflowed model ($|T|\le\auMaxProbeTwo$). On the evaluation noise of E1 the 1-rectified sampler reaches $|T|=\auMaxEvalOne$; on the
probes, for the prompts of the E2 admission grid (rotations up to $\pm45^\circ$, outside the approved range), $\auMaxGridOne$; and its
reflow pairs reached $\auMaxReflow$ before trimming. The note mentions this
instability only to justify the trimming of the reflow pairs (\lin{652--653}).
\item Between two safe prompts, probe~$\auProbe$ alone carries a median $\auShareTwoMed\%$ (mean $\auShareTwoMean\%$, 90th percentile
$\auShareTwoNinety\%$) of the squared $L^2$ map distance, and a median $\auShareFourMed\%$ of the sum that defines the $L^4$ distance.
\item The unsafe mode suppresses the blow-up: probe~$\auProbe$ leaves radius $4$ for $\auBlowSafe\%$ of the safe prompts but
$\auBlowa\%$, $\auBlowb\%$, $\auBlowc\%$ and $\auBlowd\%$ of the prompts with $m=0.01,0.03,0.1,0.3$. Hence the inverted dose-response of the
uniform $L^4$ region, which excludes $\auOneLfourUAlla$ of the requests with $m=0.01$ and $\auOneLfourUAlld$ of those with $m=0.3$: the
signature of a nuisance variable, not of a less sensitive distance.
\item Removing that one probe (Table~\ref{tab:probe}) removes every effect the note attributes to the distance: uniform $L^4$ and
$L^8$ exclude $\auOneLfourUDropd$ and $\auOneLeightUDropd$ at $m=0.3$; the nearest-field $L^2$ region excludes $\auOneLtwoNDropc$ at $m=0.1$ on the
1-rectified model, as $L^1$ and truncated $L^1$ do; in E2, nearest-field $L^2$ excludes $\auELtwoNDropc$ at $m=0.1$ and $\auELtwoNDropb$ at
$m=0.03$, against $\auETrNAllb$ for truncated $L^1$. The bounded distances barely move.
\end{itemize}
\textbf{Consequence.} The mechanism of Step~4 is wrong, and the claimed reversal of Proposition~7.3 (``$L^2$ is not the best
metric'') does not exist: once the outlier is removed, $L^2$ is at least as sensitive as truncated $L^1$. Truncation is still
useful here, but as protection against a sampler that diverges on rare noise, which the note does not discuss.
\end{findbox}

\begin{table}[!htbp]
\centering
\caption{Probe sensitivity (F1): exclusion of unsafe requests and median budget with all 256 probes $\to$ without
probe~$\auProbe$, the only probe that some pool prompt's sampler sends beyond the trimming radius~$4$. Same splits, scores and
thresholds as the note; coverage stays within $[\auProbeCovMin,\auProbeCovMax]$ in every row.}
\label{tab:probe}
\small
\begin{tabular}{l cccc c}
\toprule
& \multicolumn{4}{c}{excluded at $m=$} & \\
\cmidrule(lr){2-5}
score & $0.01$ & $0.03$ & $0.1$ & $0.3$ & budget\\
\midrule
\inputrows{tables/audit_probe.tex}
\bottomrule
\end{tabular}
\end{table}

\begin{fixbox}{Fix for F1}
State the sampler's instability where the models are introduced; replace the mechanism of Step~4 by the probe diagnosis;
justify the truncated distance by its boundedness (robustness to diverging trajectories, and a certificate with a bounded
variable), not by sensitivity; report Table~\ref{tab:probe}; correct the caption of Figure~8 and Section~10.1. A cleaner
long-term repair is a sampler that cannot diverge (clipping, or features with bounded extrapolation).
\end{fixbox}

\begin{findbox}{F2 \quad The pull-back region of the 1-rectified model is attackable \hfill \Sconc}
\textbf{What the note says.} Section~10.4 (\lin{1832--1833}): the pull-back metric ``with adaptive weights [\dots] caps the attack at a
few percent on both models''; Table~8 reports the worst attack $P(A)$ falling from $\etAtkOneEA$ (Euclidean prompt region) to
$\etAtkOnePBA$ (pull-back, 1-rectified, adaptive).

\smallskip\textbf{What the data show.} The pull-back Gram matrix $Q=U^\top G_0U$ of the 1-rectified model has eigenvalues from
$\auQminOne$ to $\auQmaxOne$ (condition number $\auQcondOne$). The note's attack proposes Gaussian steps in all 43 coordinates of the
embedding, so nearly every proposal moves along a stiff direction and is rejected. \emph{Attack~2} searches instead the five
eigen-directions of the region's metric with the smallest eigenvalues (for a field distance: of $J^\top J$, $J$ the Jacobian of
the probe outputs), from the first twelve admitted approved prompts of the same split, 400 steps each, and reports $P(A)$ on
fresh noise. It finds admitted embeddings with $P(A)=\auPBoneAtkTwoA$ (adaptive) and $\auPBoneAtkTwoS$ (static) inside the pull-back region of
the 1-rectified model (Table~\ref{tab:attack}). It does not break the field regions (worst $\auTRoneAtkTwoS$ and $\auTRoneAtkTwoA$ for the
1-rectified model, $\auTRtwoAtkTwoS$ and $\auTRtwoAtkTwoA$ for the reflowed one), nor the pull-back region of the reflowed one-step model
($\auPBtwoAtkTwoS$ and $\auPBtwoAtkTwoA$), whose pull-back metric is its field metric; on the E1 regions its worst value is $\auEoneAtkTwoMax$.

\smallskip\textbf{Why.} This is what Lemma~4.3(iii) predicts: the exact field at $t=0$ depends on the target only through its mean, so
a metric on the slice at $t=0$ cannot constrain the outputs of a curved 50-step sampler. Low-eigenvalue directions of $Q$ barely
move the slice at $t=0$ and move the later times freely.
\end{findbox}

\begin{table}[!htbp]
\centering
\caption{Worst $P(A)$ (fresh noise) found inside each region by the note's attack (12 random-search runs over all embedding
coordinates, from enforced requests with $m=0.3$) and by attack~2 (12 runs in the five directions the region penalises least,
from admitted approved prompts). The Euclidean metric has no preferred direction, so attack~2 is not defined for it.}
\label{tab:attack}
\small
\begin{tabular}{l c cc}
\toprule
& note's attack & \multicolumn{2}{c}{attack 2}\\
\cmidrule(lr){3-4}
region & worst & worst & median\\
\midrule
\inputrows{tables/audit_attack.tex}
\bottomrule
\end{tabular}
\end{table}

\begin{fixbox}{Fix for F2}
Restrict the robustness claim of the pull-back metric to the reflowed one-step model, report attack~2 for every region, and
state that a pull-back of the slice at $t=0$ of a curved model is not a safety device.
\end{fixbox}

\begin{findbox}{F3 \quad The prompt-space baseline is the weakest one available \hfill \Sconc}
\textbf{What the note says.} Abstract (\lin{108--109}) and Section~9, Step~8 (\lin{1476--1477}): ``constraining fields rather than
prompts is what protects against adversarial prompts''; Section~10.3 (\lin{1771--1773}): a prompt region ``fails, static or
adaptive, under attacks on the embedding or wording that the model ignores. A field region resists both.''

\smallskip\textbf{What the data show.} Every prompt region of the note uses the Euclidean distance between embeddings. Replace it by
the Mahalanobis distance with the covariance $S$ of the approved embeddings of the \emph{tuning} pool, regularised by
$\rho\,\mathrm{tr}(S)/43$, and tune $\rho$ and the adaptive weights on the tuning pool exactly as the note tunes its other regions
($\rho=\auMHRhoS$; adaptive weights: Gaussian kernel with $c=\auMHcAd$). On the evaluation pool:
\begin{itemize}[leftmargin=1.3em,itemsep=1pt,topsep=2pt]
\item coverage $\auMHCovS$ (static) and $\auMHCovA$ (adaptive); exclusion $\auMHExSa$, $\auMHExSb$, $\auMHExSc$, $\auMHExSd$ at $m=0.01,0.03,0.1,0.3$ (static) ---
including $m=0.01$, of which no field region excludes more than $\etDensOneF$ even with 240 anchors;
\item on the 1-rectified model, the note's attack finds at most $P(A)=\auMHAtkOneS$ (static) and $\auMHAtkOneA$ (adaptive), against
$\auFieldAtkMin$--$\auFieldAtkMax$ for the field regions; attack~2 at most $\auMHAtkTwoS$ and $\auMHAtkTwoA$;
\item its back-off is crude: requests with $m=0.3$ end at $P(A)=\auMHEnfS$ but with MMD$^2\times10^3=\auMHMmdS$ (static) and $\auMHMmdA$
(adaptive) to the request's safe target, against $\etMmdOneEA$ and $\etMmdOneFA$ for the adaptive Euclidean prompt and field regions;
\item with the nuisance token of Section~10.3, at every scale $\lambda$, exclusion at $m=0.1$ stays at least $\auMHNuMin$.
\end{itemize}
\textbf{Why.} The approved embeddings of the toy are degenerate: their third coordinate $m/0.2$ is identically $0$ (variance $\auMvar$ on
the tuning pool) and they lie on a two-dimensional surface, parametrised by rotation and radius, of $\RR^{43}$. Any detector that
whitens by the approved covariance flags every unsafe request ($m>0$) and every embedding that leaves the surface, which is what
every attack does. The comparison is therefore between field regions and the weakest prompt region, not between fields and prompts.
\end{findbox}

\begin{fixbox}{Fix for F3}
Report the Mahalanobis baseline; restrict the conclusion to Euclidean prompt distances; state that the toy cannot separate
prompt-space from field-space regions beyond that, because its approved embeddings are degenerate and its encoder exposes~$m$.
A meaningful comparison needs approved embeddings that span their space and an encoder that hides the unsafe content.
\end{fixbox}

% =====================================================================================
\section{Mathematics}\label{sec:maths}
% =====================================================================================

\begin{findbox}{M1 \quad ``One-step models make [the improved region] convex'' (\lin{1553}) \hfill \Sfalse}
The improved region is $\{\bth:\min_id_r(\bth,\bth_i)\le(M+1)\varepsilon\}=\bigcup_i\{\bth:d_r(\bth,\bth_i)\le(M+1)\varepsilon\}$, a union of
balls around the anchors, which is not convex for any sampler as soon as two balls are disjoint. Each ball is not convex either:
for a one-step sampler, $\bth\mapsto T_{\bth}(\bx)$ is affine, but $\min(\|\cdot\|,r)$ is not convex and its average over $\bx$ has
non-convex sublevel sets.
\end{findbox}
\begin{cexbox}{Counterexample}
One dimension, two noise points with equal weight, $T_\theta(x_1)=\theta$, $T_\theta(x_2)=\theta-10$, anchor $\theta_1=0$, $r=1$:
$f(\theta)=\tfrac12\min(|\theta|,1)+\tfrac12\min(|\theta-10|,1)$ gives $f(0)=f(10)=0.5$ and $f(5)=1$. The ball $\{f\le0.5\}$ contains $0$ and
$10$ but not $5$.
\end{cexbox}
\noindent The same sentence reappears in Section~11 (\lin{1880--1881}: ``one-step samplers resolve [the tension between convexity and
safety]'') and contradicts Section~10.4 (\lin{1818--1820}), which correctly calls the adaptive region non-convex. Convexity holds
for one-step samplers only with static weights and a Mahalanobis distance (Corollary~7.2(i)), the configuration the
experiments show to fail. \textbf{Fix:} delete the claim; state that the improved region is non-convex and enforced by back-off.

\begin{findbox}{M2 \quad ``A truncated map metric whose certificate needs no Markov step'' \hfill \Sfalse}
Abstract \lin{108}; ``the factor $1/r$ of Theorem~7.1(ii) is gone'' (\lin{1084--1085}); ``the absence of a Markov factor''
(\lin{1552}); ``without Markov factor'' (\lin{1871}). The proof of Proposition~7.5 rests on
$\mathbbm1\{\delta\ge r\}\le\min(\delta,r)/r$, i.e.\ on $\PP(Z\ge r)\le\E Z/r$ for $Z=\min(\|Y-Y'\|,r)$, since $\{Z\ge r\}=\{\|Y-Y'\|\ge r\}$:
Markov's inequality applied to the truncated displacement. The factor $1/r$ has moved into the normalisation of $d_r$; it is not
gone. What truncation buys is real but different: $d_r\le d_{L^1}/r$, because displacements beyond $r$ are not charged, so the
certificate is never worse than the $L^1$ one and is much better when large displacements are present. Only the tail score
$\PP(\|T_{\bth}-T_{\bth'}\|\ge r)$ avoids Markov's inequality, and it is not a metric. \textbf{Fix:} ``Markov's inequality on a bounded
variable, which does not charge displacements beyond $r$''.
\end{findbox}

\begin{findbox}{M3 \quad The size of the Gr\"onwall factor (\lin{149--150}, \lin{379--380}) \hfill \Sover}
Remark~2.14 derives $\|X_1-X_1'\|\le(\|X_0-X_0'\|+b)e^L$ and says that this factor ``is about $10^{15}$ in the toy''. The
$10^{15}$ comes from the proof-of-concept note: it is the discrete product $\prod_{l=1}^{J-1}(1+hL_l)$ with an \emph{a-priori} bound
$L_l\in[\lipmin,\lipmax]$, obtained from the triangle inequality over the features (\file{make\_tables.py}, Gr\"onwall block). With $L=\auLipMin$,
$e^L\approx10^{\auExpLLog}$; and either number bounds the factor from above, through a bound on the Lipschitz constant that may be far from
tight. The conclusion (the Gr\"onwall route is vacuous here) is plausible but not established by this number. \textbf{Fix:} ``an
a-priori bound on this factor is about $10^{15}$''.
\end{findbox}

\begin{findbox}{M4 \quad The floor of Proposition~5.6(iii) is not the mechanism of the detection failure \hfill \Sover}
Section~9, Step~5 (\lin{1350--1353}): the unsafe request's score stays below the threshold, and ``this is the floor of
Proposition~5.6(iii)''; likewise \lin{1686--1687} and \lin{1807}. The floor bounds the \emph{smallest} score, half the mean
distance between anchors. Detection depends on the margin between the threshold and the scores of typical prompts. On the
1-rectified model (100 splits) the budget $(M+1)\varepsilon$ is $\auBudgetOne$, the floor $\auFloorOne$ (ratio $2.5$), and a typical safe
test prompt has a mean distance $\auTypicalOne$ to the anchors; the margin $\auBudgetOne-\auTypicalOne$ is the spread of the safe scores,
set by rotation. For $L^2$, a far anchor also adds the unsafe displacement in quadrature,
$d(u,\bth_i)\approx(d_i^2+\delta^2)^{1/2}$, which is why averaging dilutes it. \textbf{Fix:} attribute the failure to the spread of
the safe scores; keep the floor as a statement about the tightness of the certificate.
\end{findbox}

\begin{findbox}{M5 \quad ``Of the three, the $L^2$ map metric is the most sensitive'' (\lin{1043--1048}) \hfill \Sover}
Proposition~7.3 is correct. The sentence after it compares the guaranteed lower bounds $\sqrt m\Delta$ ($L^2$), $m\Delta$ ($L^1$) and
$\log\frac1{1-m}$ (forward divergence), but a request is excluded when its score exceeds a threshold that is calibrated separately
for each distance and grows with it. Sensitivity is a ratio between the signal and the spread of the safe scores. After F1 the
claim happens to hold in the experiments, but the argument does not establish it. \textbf{Fix:} compare the lower bound with the
calibrated budget of the same distance.
\end{findbox}

\begin{findbox}{M6 \quad Minor mathematical wording \hfill \Sedit}
\begin{itemize}[leftmargin=1.3em,itemsep=1pt,topsep=2pt]
\item \lin{645}: ``fidelity can only be lost'' by reflow. Not a theorem: the exact reflowed flow reproduces $\nu_v$, and a refit error
can move the outputs towards $p$ as well as away. Here fidelity drops (MMD$^2\times10^3$ from $\eoneMmdA$ to $\eoneMmdB$).
\item \lin{655--656}: Lemma~3.9(iii) is applied to ``the exact reflowed field'' of E1 because ``its sampler map is Lipschitz''.
True, but the constant is so large (the sampler sends rare noises to $\auMaxReflow$) that the interval $t<1/(1+L)$ on which the lemma
describes the field is negligible; what is used is the identity at $t=0$, which survives. The trimming of the pairs with
$|T(\bx_0)|>4$ is harmless: the pairs are $(\mathrm{id},T)_\#\mu_0'$ for the conditioned noise law $\mu_0'$, and the proof of (iii) applies
to them verbatim. The statement should say both.
\item \lin{1389--1394}: the back-off of Step~6 is the field of a mixed prompt because the model is affine in the prompt
(Lemma~2.19(i)), not by Corollary~6.3, which concerns the projection; and the plane of Figure~9 contains affine combinations with
negative weights, not only mixtures.
\item \lin{1551}: ``this keeps [\dots] the order-$1$ Wasserstein structure''. With a Dirac weight on the candidate's nearest anchor the
reference law depends on the candidate, and the score is the nearest-anchor distance; the region is the union of balls that the note
itself compares with the order-$\infty$ region (\lin{1083--1084}). The Wasserstein reading is literal but carries no content (see R9).
\end{itemize}
\end{findbox}

% =====================================================================================
\section{Methodology and reporting}\label{sec:method}
% =====================================================================================

\begin{findbox}{R1 \quad Coverage is misreported for the tail score \hfill \Srep}
Step~5 (\lin{1359}): ``coverage lies between $\eoneCovAT$ and $\eoneCovAE$ for every distance and model''. These are the uniform scores of phase~1.
Figure~8, which the sentence cites, also shows the nearest-field scores; over the same 100 splits their coverage ranges from
$\auFigCovMin$ to $\auFigCovMax$, and the tail score reaches $\auTailCovOne$ (1-rectified) and $\auTailCovTwo$ (reflowed). The macro
\mac{eoneKnnTailCov} ($=\eoneKnnTailCov$) is computed by \file{e1\_tables.py} and never printed. The cause is ties: the tail score counts probes,
and its median threshold is $\auTailProbesOne$ probes out of 256 on the 1-rectified model and $\auTailProbesTwo$ on the reflowed one (with $r=1$
on the reflowed model the threshold is $\auTailOneEpsTwo$ and coverage $\auTailOneCovTwo$). Over-coverage is valid, but the region is then
decided by a handful of probes, and the ``tightest certificate'' of the note rests on that budget. \textbf{Fix:} report the
range of every score shown; explain the ties; randomise ties or use more probes.
\end{findbox}

\begin{findbox}{R2 \quad ``No enforced field violates its certificate on fresh noise'' cannot fail \hfill \Srep}
Table~5 and \lin{1438}. In \file{e1\_phase2.py} (\texttt{certify}, lines~96--119, and \texttt{run\_method}, line~135) the term
$q_i(A^{+r})$ of the bound and the reported $P(A)$ are both measured on the same 2000 evaluation noises. On that sample,
$\mathbbm1\{Y\in A\}\le\mathbbm1\{Y'\in A^{+r}\}+\mathbbm1\{\|Y-Y'\|\ge r\}$ holds noise by noise, so the empirical inequality can only fail
if the fresh distance exceeds the probe budget by more than the slack of the bound; for the median enforced field of the
improved method the bound is $\auSlackTrBound$ for $P(A)=\auSlackTrPA$ (tail score: $\auSlackTailBound$ for $\auSlackTailPA$). The ``0/2400'' therefore
restates that no fresh distance exceeds its budget by an order of magnitude, which the next column already reports
(at most $\eoneTwoTruncMaxRatio$ times). \textbf{Fix:} report the slack and drop the claim, or test the bound on noise independent of every
term (Section~\ref{sec:repro}).
\end{findbox}

\begin{findbox}{R3 \quad The verification step of the improved method was never run \hfill \Srep}
Step~7 (\lin{1441--1445}) proposes to verify each enforced field on $L$ fresh probes with Hoeffding's inequality, and ``The
improved method'' (\lin{1545--1551}) lists that verification as a step and concludes that ``the deployed outputs then satisfy''
the certificate. No script implements it; the deployed fields of Tables~4--5 are plug-in, and \eoneTwoPlugMin--\eoneTwoPlugMax{} of them exceed their
budget on fresh noise. The conclusion also omits the probability: with the verification it holds with probability at least
$1-\delta$. Finally, the margin matters: with $L=2000$ and $\delta=0.05$, Hoeffding adds $\auHoeff$, i.e.\ $\auHoeffTailPct\%$ of the median tail
budget ($\auBudgetTail$) and $\auHoeffTrPct\%$ of the truncated one ($\auBudgetTr$). For the tail score, a proportion, the Clopper--Pearson bound
is much tighter (no event in 2000 gives $\auCPzero$). \textbf{Fix:} implement and evaluate the verification, with Clopper--Pearson
for the tail score; until then, present it as untested.
\end{findbox}

\begin{findbox}{R4 \quad ``Field-space regions cap the unsafe mass at a few percent'' (abstract, \lin{108--109}) \hfill \Sover}
The field-space region as written admits prompts with $P(A)=\eoneTwoAttackUL$ (E1, Table~5) and $\etAtkOneFS$ (E2, static); only the localised
regions stay at a few percent ($\eoneTwoAttackLone$ and below). The evidence is twelve random-search runs on a single split, a lower bound on
the worst case; the bounds that hold for every field of the region are $\auRegionBoundTr$ (truncated) and $\auRegionBoundTail$ (tail), and they are plug-in.
\textbf{Fix:} ``localised field-space regions'', and state what the evidence is.
\end{findbox}

\begin{findbox}{R5 \quad The recommended configuration hides a safety cost (\lin{1859--1860}) \hfill \Srep}
Q3 recommends the reflowed model sampled in one step, ``at the fidelity cost of one-step sampling''. That model puts
$P(A)=\auOneStepPA$ of the outputs of \emph{approved} prompts in $A$, $\auOneStepRatio$ times the 1-rectified model ($\auRectOnePA$); this is comparable to
what the adaptive regions of the 1-rectified model leave after enforcing requests with $m=0.3$ ($\etEnfOneEA$--$\etEnfOneFA$). The number is
computed (\mac{eoneOneStepPA}) and not reported. \textbf{Fix:} report it next to the fidelity cost.
\end{findbox}

\begin{findbox}{R6 \quad Small-mass detection limits are confounded by under-generation \hfill \Sover}
The note reads the low exclusion at $m=0.01$ and $m=0.03$ as limits of the regions (\lin{1500--1501}, \lin{1671--1672},
\lin{1817--1818}). The 1-rectified model realises an excess $P(A)$ over approved prompts of $\auExcessa$ at $m=0.01$ and $\auExcessb$ at $m=0.03$,
i.e.\ $\auRealiseda\%$ and $\auRealisedb\%$ of the targets' $\auTargeta$ and $\auTargetb$: these fields are much closer to safe fields than their
nominal mass suggests. \textbf{Fix:} report exclusion against the realised excess $P(A)$, or state the realised masses.
\end{findbox}

\begin{findbox}{R7 \quad Statistical precision \hfill \Srep}
No interval accompanies any number of Sections~9--10. Exclusion at a mass is a mean over 100 distinct requests (E1 and E2), so a
binomial 95\% half-width is up to $\pm0.10$ whatever the number of splits; enforcement uses 5 (E1) and 3 (E2) splits; every attack
uses one split and twelve runs. In E1, the improved variants (nearest field, $r=1$, tail at $0.5$) were selected by a screen on the
same pool as the numbers reported for them; E2 repairs this with disjoint tuning and evaluation pools and confirms the selection.
\textbf{Fix:} state the precision and the selection.
\end{findbox}

\begin{findbox}{R8 \quad External validity of nearest-anchor detection \hfill \Sscope}
The approved family is two-dimensional (rotation, radius). The distance from a request to its nearest anchor shrinks like
$M^{-1/d}$ for a $d$-dimensional family, and already here $m=0.03$ needs 240 anchors (Figure~16). Nothing in the note discusses
what happens when the approved prompts span many dimensions. \textbf{Fix:} state the dependence on the intrinsic dimension.
\end{findbox}

\begin{findbox}{R9 \quad Framing and related work \hfill \Sscope}
With weights on the nearest anchor, the region is the split-conformal nearest-neighbour novelty detector
(Laxhammar and Falkman, \emph{Ann.\ Math.\ Artif.\ Intell.}\ 74, 2015; Bates et al., \emph{Ann.\ Statist.}\ 51, 2023), applied to
a distance between sampler maps; the experiments favour it over the order-$1$ score that gives the note its title. Prompt-space
regions with a whitened covariance are a standard out-of-distribution detector (Lee et al., NeurIPS 2018). None is cited.
\textbf{Fix:} say so where the improved method is stated, and cite.
\end{findbox}

\begin{findbox}{R10 \quad ``Tested: [\dots] prompt mixing'' (\lin{1872--1874}) \hfill \Srep}
What E1 checks is the identity of Corollary~6.3 (largest deviation $\auMixcheck$ in \file{eval\_enforce.json}), which holds by
linearity. The projection of Proposition~6.2(iii), the constructive result of Section~6, was computed in phase~1
(\file{e1\_eval.py}) and its effect on $P(A)$ is never reported; every enforcement in the note is a back-off. \textbf{Fix:} say what was
checked, or report the projection.
\end{findbox}

% =====================================================================================
\section{Editorial and bookkeeping}\label{sec:edit}
% =====================================================================================
\begingroup\sloppy
\begin{itemize}[leftmargin=1.6em,itemsep=2pt,topsep=2pt]
\item[H1] \file{README.md} calls the note 23 pages; it has 29.
\item[H2] \file{e2\_figures.py}, lines~325--327: the spread \mac{etSpreadOne} averages over the evaluation prompts against \emph{all}
1000 approved prompts, diagonal zeros included, and the displacement \mac{etDispOne} uses the 200 bases of both pools and the
nearest of all 1000 approved prompts, whereas the text describes evaluation-pool quantities. The values barely change
(spread $\auSpreadOld$ against $\auSpreadNew$ between distinct evaluation prompts; displacement $\auDispOld$ against $\auDispNew$).
\item[H3] Q1 (\lin{1854--1856}) speaks of weights ``that take the transport mass from the nearest anchors''; the adaptive field
region of the reflowed model, tuned on the tuning pool, uses prompt-kernel weights (Table~7: ``$c=0.3$, prompt'').
\item[H4] Table~2: the rows ``$G_{\mathrm{on}}$, ODE'' and ``reverse KL, SDE'' agree to two decimals in all five columns. This is a
coincidence of rounding, not a copy error: at $m=0.01$ and $m=0.03$ the exclusions are $\auPocOnOne$, $\auPocOnThree$ and $\auPocRevOne$, $\auPocRevThree$.
A caption note avoids the question.
\item[H5] \lin{11}: ``Every number in Section~7 is copied from \file{tables/numbers.tex}'' --- that section is now Section~8, and its
Table~2 is typed by hand.
\item[H6] Computed but never printed: \mac{eoneKnnTailCov}, \mac{eonePlugFail}, \mac{eonePlugMax}, \mac{eoneRegionWorst}, \mac{eoneOneStepPA},
\mac{etTuneGap}, \mac{etPoolPATwo}; R1 and R5 show that two of them change the reading of the results.
\item[H7] Section~8: the Clopper--Pearson bound on $\bar P_w(A)$ (\file{part\_d.py}, line~82) treats the $M\times4000$ indicators as
independent, but the $M$ anchors share the same 4000 noise chains; the bound is slightly anti-conservative.
\end{itemize}
\endgroup

% =====================================================================================
\section{What was checked and holds}\label{sec:ok}
% =====================================================================================
\begingroup\sloppy
\noindent\Ok{} \textbf{Section~2.} Definitions~2.1, 2.3, 2.7, 2.9, 2.11, 2.15, 2.17, 2.20, 2.21; Lemma~2.4 (including (iii): truncation
of a metric is a metric, and (iv): integrals of pseudometrics); Lemma~2.5; Theorem~2.6 (i)--(iii); Proposition~2.8; Lemma~2.10,
including Bayes' formula with the null set where the ratio is undefined; Lemma~2.13; the Gr\"onwall inequality of Remark~2.14;
Lemma~2.16; Lemma~2.19.

\noindent\Ok{} \textbf{Section~3.} Remark~3.2, including ``for almost every $s$'' from conditioning on $(\bx_t,t)$ jointly;
Lemmas~3.3 and~3.6; Lemma~3.9: (ii) through the dyadic-cube argument and two versions of $\E[\mathbbm1_B(\bx_1)\mid\bx_t,t]$, (iii)
through the contraction $\bx\mapsto(\by-tT(\bx))/(1-t)$, (iv) through strong monotonicity of $\Psi_t$ and the induction on the
Euler iterates; the identity of Liu et al.\ behind Remark~3.10, re-derived (crossing plus curvature) and checked against the
figure's numbers ($\rfDecrease$ against $\rfSum$).

\noindent\Ok{} \textbf{Sections~4--7.} Proposition~4.2; Lemma~4.3 (the Bayes version is continuous on $\RR^n\times[0,1)$ and equals
$\E_p[\bx_1]-\bx$ at $t=0$); Lemma~5.3 (plan and dual certificate); Proposition~5.4 (the rank argument holds conditionally on $\cD$
in the sense of Definition~2.21); Proposition~5.6; Proposition~6.2, including the convex-hull argument of (iii); Corollary~6.3;
Proposition~6.4; Theorem~7.1; Corollary~7.2; Proposition~7.3; Proposition~7.5 (the inequality
$\mathbbm1\{\delta\ge r'\}\le\max(1,r/r')\min(\delta,r)/r$ in both cases $r'\le r$ and $r'>r$).

\noindent\Ok{} \textbf{Appendix~A.} The drift and diffusion of Definition~A.1 match the marginal-preserving SDE of the linear
interpolant; Lemma~A.2 (chain rule over Gaussian kernels); Proposition~A.3 (i)--(iii).

\noindent\Ok{} \textbf{Experiments.} Split-conformal thresholds implement $C_{1-\beta}$ exactly (\texttt{kth\_threshold}); every
reported coverage is within sampling error of $55/61$ or $37/41$ (or above it, R1); tuning and evaluation pools of E2 are disjoint, and
the tuned weights are within $\etTuneRegret$ of the best on the evaluation pool; the pull-back metric of the reflowed model equals its
one-step map metric to $7\times10^{-8}$; the attack reports $P(A)$ on noise not used by the search; static regions are too wide on
this family (E1 and E2), and nearest-anchor weights repair detection; field regions resist both attacks (Table~\ref{tab:attack}).

\endgroup

% =====================================================================================
\section{Resolution in the revised note}\label{sec:resolution}
% =====================================================================================
The note was revised on 14 October 2026, in place (\file{safe\_prompt\_fields.tex}, now 33 pages); the audited version is archived.
Section, table and figure numbers below refer to the revised note.

\renewcommand{\arraystretch}{1.15}
\begin{longtable}{@{}p{1.0cm}>{\raggedright\arraybackslash}p{15.2cm}@{}}
\toprule
\textbf{Id} & \textbf{Resolution}\\
\midrule
\endhead
F1 & The sampler's divergence is stated where the models are introduced (Section~9, Step~1; Remark~3.10). Step~4 replaces the ``blob
changes'' mechanism by the probe diagnosis, with a table of every score with and without the flagged probe (Table~4); Step~5 and the caption of
Figure~8 (which now shows nearest-field $L^2$ without the probe) give $L^2$ its corrected standing; Section~10.1 justifies the truncated
distance by its boundedness. The paragraph after Proposition~7.3 warns that unbounded distances inherit far-away trajectories.\\
F2 & The directed search (attack~2) is part of both experiments: Section~9, Step~8, and its certificate table (Table~7), Table~9 and
Figure~14 of Section~10, and the pull-back rows of Table~10. Section~10.4 states that the pull-back metric of the 1-rectified model is
not a safety device, with the directed-search numbers; the summary (Q3) and Section~11 say so too.\\
F3 & The Mahalanobis prompt region is defined in Section~10.1 (Table~8), evaluated in Table~9, Figures~14--15 and Table~10, and
discussed in a new paragraph of Section~10.3, including its crude enforcement. The abstract, Step~8, the take-away of Section~10.3,
Q2 and Section~11 restrict the conclusion to Euclidean prompt distances and state why the toy cannot separate the two spaces.\\
M1 & The convexity claim is deleted (``The improved method'', Section~11); after Proposition~7.5 the union-of-balls region is stated
to be non-convex for any sampler.\\
M2 & Every ``no Markov step'' is replaced by ``Markov's inequality on a bounded variable'', and the real gain, $d_r\le d_{L^1}/r$, is
stated after Proposition~7.5.\\
M3 & Remark~2.14 and Section~1 speak of an a-priori bound on the factor, quoting the proof of concept's own macros.\\
M4 & Section~9, Step~5, and Sections~10.2 and~10.4 attribute the failure of static regions to the spread of the safe scores, with
the budget, the typical score and the floor side by side.\\
M5 & The paragraph after Proposition~7.3 compares the lower bounds with the calibrated budget of each distance.\\
M6 & Remark~3.10 (fidelity; Lemma~3.9(iii) for the trimmed pairs and a large Lipschitz constant); Step~6 (Lemma~2.19(i) for the
back-off, affine combinations in the plane of Figure~9); the ``order-$1$ structure'' sentence is gone.\\
R1 & Step~5 gives the coverage of every score of Figure~8, the over-coverage of the tail score and its cause.\\
R2 & The certificate table of Section~9 (Table~7) reports the median $P(A)$ next to the median bound instead of the violation count, and
Step~7 explains why the count was uninformative.\\
R3 & The verification is implemented (\file{code/e1\_verify.py}) and reported in Step~7 with its own table (Table~5): a two-stage check on
$2000$ fresh noises, Hoeffding margins for the truncated distance and Clopper--Pearson for the tail score, fallback to the nearest
anchor, and an evaluation on independent noise. On the 1-rectified model, the deployed fields of the improved region exceed their
budget on $10{,}000$ independent noises in $\auVerOverOneTr$ of cases (against \eoneTwoPlugMin--\eoneTwoPlugMax{} on the evaluation noise
before), and none of the $\auVerN$ deployed fields has $P(A)$ above its certificate on noise independent of every term
($\auVerViol$ violations). ``The improved method'' carries the probability qualifier.\\
R4 & The abstract speaks of localised field-space regions and two attacks.\\
R5 & Section~10.4 and Q3 report the $P(A)$ of approved prompts under one-step sampling.\\
R6 & Section~9, Step~1, gives the realised unsafe masses; the summaries of Sections~9 and~10.4 refer to them.\\
R7 & Section~9, Step~1, and Section~10.1 state the precision of exclusion rates and that every attack uses one split.\\
R8 & Section~10.4 and Section~11 state the dependence of nearest-anchor detection on the dimension of the approved family.\\
R9 & After Proposition~7.5 and in ``The improved method'': the nearest-neighbour conformal novelty detector, with Laxhammar and Falkman
(2015) and Bates et al.\ (2023); Lee et al.\ (2018) for the Mahalanobis baseline (\file{refs.bib}, Part~5, checked against the sources).\\
R10 & Step~6 says what was checked (the identity of Corollary~6.3) and that every score is enforced by back-off; Section~11 lists
``the identity of prompt mixing''.\\
H1--H7 & README page count and pipeline; \file{e2\_figures.py} computes spread and displacement on the evaluation pool; Q1 names the
prompt-kernel weights of the reflowed field region; the caption of Table~2 explains the coincidence (three-decimal macros from
\file{make\_tables.py}) and the shared-noise caveat; the header comment is corrected and Section~8 quotes the proof of concept's
macros; the two unused macros that mattered (R1, R5) are printed.\\
\bottomrule
\end{longtable}

% =====================================================================================
\section{Reproduction}\label{sec:repro}
% =====================================================================================
From \file{code/}, after the pipelines of the README (E1 and E2):
\begin{center}\small
\begin{tabular}{l>{\raggedright\arraybackslash}p{9.2cm}l}
\toprule
script & content & time\\
\midrule
\file{python e1\_robustness.py} & F1 on E1 (screen with and without probe $\auProbe$), attack~2 on the E1 regions & 5 min\\
\file{python e2\_robustness.py} & F1 on E2, attack~2 on the E2 regions, the Mahalanobis regions (F3) & 15 min\\
\file{python e1\_verify.py} & R2--R3: verified deployment, tested on independent noise & 45 min\\
\file{python audit\_checks.py} & every number and both tables of this audit & 30 s\\
\bottomrule
\end{tabular}
\end{center}
The audited note, its number macros and the figure scripts changed by the revision are archived in
\file{archive/2026-10-14\_pre\_audit/}.

\end{document}
